\subsection{Component costs}
\label{sec:costs}

Every cost in the model comes from a price that someone paid or quoted, or is marked as our own estimate. \Cref{sec:inputs} lists those that matter for the results, the fuel and carbon prices among them, with their ranges and sources, and \cref{fig:capital} shows what they add up to for a few ships.

The ship itself dominates. A new \num{3800}-tonne general cargo ship from an Indian yard costs about USD~12 million, and the price grows with deadweight to the power 0.40 up to Chinese-built bulk carriers of \num{82000}~t. The fit uses an order of six coasters at Udupi Cochin in 2023 and broker assessments of Handysize, Ultramax and Kamsarmax newbuilds in 2026; the yard country changes with size, which the fit cannot separate. A dual-fuel ammonia engine costs \SI{190}{USD/kW} more than a diesel \citep{lrumas2020}.

Every rig is priced at one production stage: a series of 50 or more units with no learning beyond it, the basis of the cost model of \citet{joshi2025}. No rig has reached it, since about 23 rotor, 26 wing sail and 6 kite installations had been made or ordered by 2023 \citep{emsa2023}. A rig's manufacturing cost is its materials and bought-in parts at today's purchase prices and its labour at USD~26 an hour, and its price is that cost times a market factor of 2 for the maker's overhead and margin. Every source figure under the factor is a cost before that overhead and margin, and a bought-in part enters at its purchase price, which holds only its supplier's margin. Maintenance costs 2\,\% of the price a year \citep{emsa2023}, the share \citet{tillig2020b} also take for rotors. A kite costs its hardware once and its canopy and tether per flight hour. The winch only reels the kite out and in at launch and landing and otherwise holds the tether on its brake, so it is a mooring winch, not the generator winch of a pumping kite power system. The launch mast is telescopic, as on the SkySails ships \citep{wrage2018}, and only holds the kite at launch and landing; the tether runs to a tow point on deck. It is 1.3 times the square root of the projected kite area tall: the \SI{15}{m} mast of the SkySails kite of \SI{125}{m^2} projected gives 1.34 \tbd{cite SkySails mast height}, and the \SI{35}{m} mast of the \SI{1000}{m^2} Airseas kite \citep{airseas_seawing_web}, whose area basis is not stated, gives 1.11 if projected and 1.25 if flat. It is at most \SI{35}{m}; its steel mass is \SI{0.5}{t} plus a part that grows with kite area to the power $2/3$ times height to the power $5/3$. Kites larger than about \SI{720}{m^2} would need a taller mast than any built and are flown here from a \SI{35}{m} mast, an extrapolation. At a fixed aspect ratio the span is a constant times the square root of the area, so this rule also is one on span. With the control unit and 15\,\% for installation, an \SI{800}{m^2} kite system costs about USD~\num{410000}, and its ram-air canopy another USD~\num{190000} every \num{4000} flight hours. We build the prices of wings and rotors up from the masses of current designs at the material costs and labour hours of the wind-turbine industry: a wing sail costs USD~\num{132000} for its controls, hydraulics and assembly plus \SI{1330}{USD/m^2} of sail, and a rotor USD~\num{148000} plus \SI{4080}{USD/m^2} of height times diameter, twice their manufacturing cost. All three rigs add 15\,\% for installation on a newbuild \citep{emsa2023}, and retractable wings and rotors add 12\,\% for the tilting, folding or rails that lower them (\cref{sec:model}). Today's market prices are the upper bounds.

Five inputs need a word about their stage, and we give each with what it does. The composite of wings and rotors is priced at the material costs of a wind blade plant making about a thousand blades a year, with its labour hours raised by a third for small runs, which may put the deck rigs nearer maturity than the kite. Their prices bound this: a rotor lands within 13\,\% below and 6\,\% above Norsepower's 2019 prices of EUR~\num{300000} to \num{750000} per rotor \citep{tillig2020b}, or up to 22\,\% above them if those include installation, as their wording allows, so the rotors stand at about the stage of the rotors sold so far. The kite hardware, built up the same way, comes to a fifth to a half of what ship kites have sold for, all of them prototypes; \cref{sec:disc_prices} gives what pricing it near those does. The canopy price goal is a series of 200 systems a year, past this stage, while today's canopy price, from retail cloth and paraglider sewing times, is on it or above it; \cref{sec:res_missed} gives what the goal is worth. The tether is rope bought from the rope maker at EUR~80/kg \citep{joshi2025}, with no factor; the factor would add 15\,\% to the wear per flight hour of a \SI{1600}{m^2} kite. The avionics of the VTOL kite, EUR~\num{150000}, are the early commercial figure of \citet{joshi2025}, which aviation certification and redundancy raise above the EUR~\num{15000} of a prototype, so the prototype figure, our low bound, is a kite flown without that certification and not an earlier stage.

The cost per tonne-kilometre stays the measure for all ships, so that classes compare like for like. For volume-limited ships it is a cost per slot of cargo work in disguise: their tonnes are their slots times a fixed cargo per unit. We therefore also report the cost per TEU-kilometre and per car-kilometre of the container ships and the car carrier, \tbd{X} and \tbd{Y} US cents, as a check.

\begin{figure}[tbp]
\centering
\plot{capital}
\caption{Capital of an \SI{85}{m} and a \SI{170}{m} ship by component, on VLSFO and on e-ammonia without a rig and on e-ammonia with a \SI{1600}{m^2} kite, \SI{800}{m^2} of wing sails or \SI{400}{m^2} of rotors: the hull and its basic drive, the extra cost of a dual-fuel ammonia engine, and the rig; the fuel tank, sized by the route, is not shown.}
\label{fig:capital}
\end{figure}
